\chapter{Zdania i spójniki logiczne}\label{ch:propositional}

\section{Po co matematyce logika?}\label{sec:logic-motivation}
Wyobraźmy sobie rozumowanie: „Jeżeli liczba jest podzielna przez $4$,
to jest parzysta. Liczba $12$ jest podzielna przez $4$, a zatem jest
parzysta”. Nie trzeba jeszcze znać żadnej zaawansowanej teorii, aby
zauważyć, że wniosek wynika z podanych przesłanek. Logika bada
\emph{formę} takich wnioskowań: co sprawia, że przejście od przesłanek
do wniosku jest poprawne, niezależnie od tego, czy rozmawiamy o
liczbach, figurach czy programach komputerowych.

Od początku oddzielamy dwa zagadnienia. \emph{Prawdziwość zdania}
dotyczy jego wartości logicznej w ustalonej interpretacji. \emph{Poprawność wnioskowania} dotyczy tego, czy
możliwa jest sytuacja, w której wszystkie przesłanki są prawdziwe,
ale wniosek fałszywy. Zobaczymy, że nie są to te same pytania.

\begin{przyklad}[Prawdziwa implikacja z fałszywym poprzednikiem]
„Jeżeli $2$ jest liczbą nieparzystą, to $2+2=4$”. Cała wypowiedź
jest prawdziwa w logice klasycznej, mimo że jej poprzednik jest
fałszywy. Znaczenie tego pozornie dziwnego faktu wyjaśnimy przy
implikacji, zamiast przyjmować go bez komentarza.
\end{przyklad}

\section{Zdania, pytania i wypowiedzi ze zmiennymi}\label{sec:statements}
\begin{definicja}[Zdanie w sensie logiki klasycznej]\label{def:statement}
\emph{Zdaniem} nazywamy wypowiedź oznajmującą, której w ustalonej
interpretacji można przypisać dokładnie jedną z dwóch wartości:
\emph{prawdę} (oznaczaną $1$) albo \emph{fałsz} (oznaczany $0$).
W tym rozdziale posługujemy się klasyczną logiką dwuwartościową.
\end{definicja}

\begin{uwaga}[Granice wstępnej definicji zdania]\label{rem:statement-scope-xli}
Samo oznajmujące brzmienie wypowiedzi nie gwarantuje, że jest
ona zdaniem o jednoznacznej wartości logicznej. W szczególności
samoodnosząca się wypowiedź „To zdanie jest fałszywe” prowadzi
przy naiwnym zastosowaniu reguły dwuwartościowości do sprzeczności:
jeżeli jest prawdziwe, to --- zgodnie ze swą treścią --- jest
fałszywe, a jeżeli jest fałszywe, to jego treść jest prawdziwa.
Dlatego w ścisłym rachunku logicznym ustalamy \emph{język}
i jego \emph{interpretację}; nie przyznajemy automatycznie wartości
logicznych wszystkim napisom języka naturalnego. Składnię
rachunku zdań zdefiniujemy niżej, a składnię logiki pierwszego
rzędu --- w rozdziale~\ref{ch:syntax-fol}.
\end{uwaga}

\begin{przyklad}[Zdania i wypowiedzi niebędące zdaniami]
Wypowiedzi „$2+3=5$” oraz „$7<4$” są zdaniami: pierwsze jest
prawdziwe, drugie fałszywe. Pytanie „Ile wynosi $2+3$?” i polecenie
„Oblicz tę sumę” nie są zdaniami w powyższym sensie. Napis „$x>4$”
nie jest jeszcze zdaniem o ustalonej wartości logicznej: jego
prawdziwość zależy od wartości $x$. Po podstawieniu $x=7$ otrzymujemy
zdanie prawdziwe, a po podstawieniu $x=2$ --- fałszywe.
\end{przyklad}

\begin{uwaga}[Zdanie a umiejętność rozstrzygnięcia prawdy]
Definicja nie wymaga, abyśmy potrafili rozpoznać, czy dane zdanie
jest prawdziwe. Matematyczne twierdzenie może pozostawać nierozstrzygnięte
przez wiele lat, a mimo to w ustalonej interpretacji klasycznej
jest zdaniem. Nie wolno utożsamiać „nie znam odpowiedzi” z
„wypowiedź nie ma wartości logicznej”.
\end{uwaga}

\begin{definicja}[Zdania proste i złożone]
Litery zdaniowe $p,q,r,\ldots$ są symbolami, które po ustaleniu
interpretacji mogą reprezentować całe zdania (na przykład „$6$ jest
parzyste”), a nie liczby. Za ich pomocą i przy użyciu spójników
budujemy \emph{formuły zdaniowe}. W rachunku zdań litery traktujemy
jako symbole formalne: w różnych wartościowaniach ta sama litera
może otrzymać różną wartość logiczną. Po zastąpieniu liter
konkretnymi zdaniami o ustalonych wartościach logicznych
otrzymujemy zdanie; jest ono złożone, jeżeli użyliśmy
spójników. Nawiasy określają budowę formuły.
\end{definicja}

\begin{definicja}[Skończona budowa formuł rachunku zdań]
\label{def:propositional-grammar-xli}
Wybierzmy litery zdaniowe $p_1,p_2,\ldots$.
\emph{Formuły rachunku zdań} powstają według następujących reguł:
\begin{enumerate}[label=(\roman*)]
\item każda litera zdaniowa jest formułą;
\item jeżeli $\alpha$ jest formułą, to $\neg\alpha$ jest formułą;
\item jeżeli $\alpha,\beta$ są formułami, to
      $(\alpha\land\beta)$, $(\alpha\lor\beta)$,
      $(\alpha\Rightarrow\beta)$ i
      $(\alpha\Leftrightarrow\beta)$ są formułami;
\item formułami są tylko napisy otrzymane przez skończoną liczbę
      zastosowań reguł (i)--(iii).
\end{enumerate}
Jest to \emph{definicja indukcyjna składni}: mówimy, z jakich
prostych napisów buduje się poprawne napisy, jeszcze bez pytania
o ich prawdziwość. Umowy dotyczące opuszczania nawiasów
podamy w sekcji~\ref{sec:biconditional}.
\end{definicja}

\begin{przyklad}[Jednoznaczny rozbiór nawiasów]\label{ex:parse-tree-xli}
Napis $p\Rightarrow(\neg q\lor r)$ ma w pełnym zapisie postać
$(p\Rightarrow(\neg q\lor r))$. Najpierw powstają
formuły $p,q,r$; następnie $\neg q$ i
$(\neg q\lor r)$, a na końcu cała implikacja.
Inny napis $((p\Rightarrow\neg q)\land r)$ ma inną budowę.
Gdy $p=0$, $q=1$, $r=0$, pierwszy ma wartość $1$
(fałszywy poprzednik implikacji), a drugi $0$
(jeden z członów koniunkcji jest fałszywy). Nie wolno więc
usuwać nawiasów w sposób zmieniający budowę formuły.
\end{przyklad}

\begin{center}
\begin{tikzpicture}[
  level distance=11mm,
  level 1/.style={sibling distance=36mm},
  level 2/.style={sibling distance=24mm},
  every node/.style={draw,rounded corners=2pt,inner sep=3pt,font=\small}]
\node {$\Rightarrow$}
 child {node {$p$}}
 child {node {$\lor$}
   child {node {$\neg$} child {node {$q$}}}
   child {node {$r$}}};
\end{tikzpicture}
\par\smallskip
\small Schemat budowy formuły
$p\Rightarrow(\neg q\lor r)$: liście są literami, węzły spójnikami.
\end{center}

\section{Negacja, koniunkcja i alternatywa}\label{sec:connectives-first}
\begin{definicja}[Trzy podstawowe spójniki]\label{def:not-and-or}
Niech $p,q$ oznaczają zdania.
\begin{itemize}
\item \emph{Negacja} $\neg p$ (czytamy: „nie $p$”) jest prawdziwa
      dokładnie wtedy, gdy $p$ jest fałszywe.
\item \emph{Koniunkcja} $p\land q$ („$p$ i $q$”) jest prawdziwa
      dokładnie wtedy, gdy \emph{oba} zdania są prawdziwe.
\item \emph{Alternatywa} $p\lor q$ („$p$ lub $q$”) jest prawdziwa
      dokładnie wtedy, gdy \emph{co najmniej jedno} z dwóch zdań
      jest prawdziwe. W szczególności oba mogą być prawdziwe.
\end{itemize}
Reguły te wyraża następująca tablica prawdy:
\[
\begin{array}{c c|c c c}
p&q&\neg p&p\land q&p\lor q\\\hline
1&1&0&1&1\\
1&0&0&0&1\\
0&1&1&0&1\\
0&0&1&0&0
\end{array}
\]
\end{definicja}

\begin{przyklad}[„Lub” w języku naturalnym]
Niech $p$ oznacza „$6$ jest parzyste”, a $q$ „$6$ jest
podzielne przez $3$”. Oba zdania są prawdziwe, a więc także
$p\lor q$ jest prawdziwe. Jest to \emph{lub niewykluczające}.
W języku codziennym „albo herbata, albo kawa” może wykluczać
jednoczesne wystąpienie obu zdarzeń. W logice matematycznej
\emph{wyłączność} trzeba wyrazić osobno: $(p\lor q)\land\neg(p\land q)$.
\end{przyklad}

\begin{przyklad}[Negacja zdania złożonego]
Ustalmy liczbę całkowitą $n$. Niech $p$ oznacza zdanie „$n$ jest
liczbą dodatnią”, a $q$ zdanie „$n$ jest liczbą parzystą”. Negacja
zdania „$n$ jest liczbą dodatnią i parzystą” nie brzmi „$n$ jest
liczbą niedodatnią i nieparzystą”. Wystarczy, że zawiedzie
\emph{jeden} z dwóch warunków: $n$ nie jest liczbą dodatnią
\emph{lub} nie jest liczbą parzystą. Za chwilę uzasadnimy to formalnie.
\end{przyklad}


\begin{uwaga}[Jak czytać nawiasy w zapisie logicznym?]
Jeżeli $p$ oznacza „pada deszcz”, a $q$ „wieje wiatr”, to
$\neg(p\land q)$ oznacza zaprzeczenie całego zdania „pada i wieje”.
Natomiast $(\neg p\land q)$ oznacza „nie pada, ale wieje”.
To nie są te same wypowiedzi: gdy pada deszcz, lecz nie wieje,
pierwsza jest prawdziwa, a druga fałszywa. Dlatego w tym skrypcie
nawiasy nie są ozdobą: wskazują dokładnie, do jakiej części
wypowiedzi należy dany spójnik.
\end{uwaga}

\section{Implikacja: kiedy „jeżeli..., to...” jest prawdziwe?}
\label{sec:implication}
\begin{definicja}[Implikacja materialna]\label{def:implication}
Dla zdań $p,q$ \emph{implikacją} $p\Rightarrow q$ nazywamy
zdanie fałszywe dokładnie wtedy, gdy $p$ jest prawdziwe, a $q$
fałszywe. Zdanie $p$ to \emph{poprzednik} (założenie), a $q$ to
\emph{następnik} (konkluzja) implikacji.
\[
\begin{array}{c c|c}
p&q&p\Rightarrow q\\\hline
1&1&1\\
1&0&0\\
0&1&1\\
0&0&1
\end{array}
\]
\end{definicja}

\paragraph{Dlaczego dokładnie taka tablica?}
Warunek „jeżeli $p$, to $q$” ma wykluczać sytuację, gdy $p$
zachodzi, a $q$ nie zachodzi. Właśnie taka sytuacja jest
\emph{kontrprzykładem} do implikacji. Gdy $p$ nie zachodzi,
obietnica nie została naruszona, niezależnie od wartości $q$.
Tak rozumiana implikacja jest regułą o wartości logicznej, a nie
stwierdzeniem związku przyczynowego między dwiema wypowiedziami.

\begin{przyklad}[Fałszywy poprzednik nie daje kontrprzykładu]
Zdanie „jeżeli $3$ jest parzyste, to $5$ jest większe od $2$”
jest prawdziwe. Prawdziwe jest także zdanie „jeżeli $3$ jest
parzyste, to $5$ jest mniejsze od $2$”. W obu przypadkach
poprzednik jest fałszywy, a zatem nie występuje układ
„prawdziwy poprzednik --- fałszywy następnik”.
\end{przyklad}

\begin{definicja}[Warunek konieczny i wystarczający]
Niech $p,q$ będą formułami zdaniowymi i niech zostanie ustalony
zbiór dopuszczalnych wartościowań ich liter. Jeżeli
$p\Rightarrow q$ jest prawdziwa przy każdym z tych wartościowań,
to $p$ nazywamy \emph{warunkiem wystarczającym} dla $q$,
a $q$ --- \emph{warunkiem koniecznym} dla $p$.
Przy braku dodatkowych ograniczeń rozważamy wszystkie
wartościowania. Nie wystarcza sprawdzenie prawdziwości
implikacji tylko przy jednym wartościowaniu, na przykład takim,
w którym $p$ jest fałszywe. Jeżeli przy każdym dopuszczalnym
wartościowaniu prawdziwe są obie implikacje $p\Rightarrow q$
i $q\Rightarrow p$, to $p$ jest dla $q$ warunkiem zarówno
koniecznym, jak i wystarczającym. Są to sposoby opisu
zależności logicznych, a nie osobne twierdzenia.
\end{definicja}

\begin{przyklad}[Konieczny nie znaczy wystarczający]
Rozpatrujemy liczby całkowite $n$ oraz zdania „$4\mid n$” i „$2\mid n$”.
Dopuszczalne są tutaj tylko te pary wartości logicznych, które
pochodzą od tej samej liczby $n$, a nie wszystkie niezależne
przypisania wartości dwóm literom zdaniowym. Jeżeli $4\mid n$,
to $n=4k=2(2k)$ dla pewnego $k\in\mathbb Z$, zatem $2\mid n$.
Podzielność przez $4$ jest więc warunkiem wystarczającym
parzystości, a parzystość --- koniecznym warunkiem podzielności
przez $4$. Nie jest jednak warunkiem wystarczającym: liczba $6$
jest parzysta, ale nie jest podzielna przez $4$. Nie oznacza to,
że sama formuła $p\Rightarrow q$ jest tautologią, gdy $p$ i $q$
traktujemy jako niezależne litery zdaniowe.
\end{przyklad}

\begin{definicja}[Implikacja odwrotna, przeciwna i przeciwstawna]
Dla $p\Rightarrow q$ rozważamy:
\begin{itemize}
\item \emph{implikację odwrotną} $q\Rightarrow p$;
\item \emph{implikację przeciwną} $\neg p\Rightarrow\neg q$;
\item \emph{implikację przeciwstawną} $\neg q\Rightarrow\neg p$.
\end{itemize}
W dalszym ciągu udowodnimy, że implikacja jest logicznie
równoważna swojej implikacji przeciwstawnej. Implikacje odwrotna
i przeciwna na ogół nie są z nią logicznie równoważne.
\end{definicja}

\begin{przyklad}[Dlaczego nie wystarczy odwrócić strzałki?]
„Jeżeli liczba jest podzielna przez $4$, to jest parzysta”
jest prawdziwe. Zdanie odwrotne „jeżeli liczba jest parzysta,
to jest podzielna przez $4$” jest fałszywe: kontrprzykładem
jest $6$. Zdanie przeciwstawne „jeżeli liczba nie jest
parzysta, to nie jest podzielna przez $4$” jest prawdziwe.
\end{przyklad}

\section{Równoważność i inne użyteczne skróty}\label{sec:biconditional}
\begin{definicja}[Równoważność]\label{def:biconditional}
Dla zdań $p,q$ formuła $p\Leftrightarrow q$
(czytamy: „$p$ wtedy i tylko wtedy, gdy $q$”)
jest prawdziwa dokładnie wtedy, gdy $p$ i $q$ mają tę samą
wartość logiczną. Odpowiada koniunkcji obu implikacji:
$(p\Rightarrow q)\land(q\Rightarrow p)$.
\[
\begin{array}{c c|c}
p&q&p\Leftrightarrow q\\\hline
1&1&1\\
1&0&0\\
0&1&0\\
0&0&1
\end{array}
\]
\end{definicja}

\begin{przyklad}[Dowód równoważności wymaga dwóch kierunków]
Zdanie „liczba całkowita $n$ jest parzysta wtedy i tylko wtedy,
gdy $n$ da się zapisać jako $2k$ dla pewnej liczby całkowitej $k$”
to w istocie definicja parzystości. Aby udowodnić \emph{inną}
równoważność, często trzeba odrębnie przeprowadzić dowody
$p\Rightarrow q$ oraz $q\Rightarrow p$.
\end{przyklad}

\begin{uwaga}[Umowa dotycząca nawiasów]
Najwyższy priorytet ma negacja $\neg$, potem koniunkcja
$\land$, następnie alternatywa $\lor$, implikacja $\Rightarrow$
i równoważność $\Leftrightarrow$. Jeżeli jakiekolwiek wątpliwości
są możliwe, zapisujemy nawiasy. Dla przykładu
$\neg p\land q$ oznacza $(\neg p\land q)$, a nie
$\neg(p\land q)$. Ciągi spójników dwuargumentowych tego samego
priorytetu również zapisujemy z nawiasami: nie posługujemy się
niejednoznacznym zapisem $p\Rightarrow q\Rightarrow r$.
Skrót $\varphi_1\land\cdots\land\varphi_n$ ($n\ge1$)
oznacza natomiast koniunkcję grupowaną od lewej: jest to
$C_n$, gdzie $C_1=\varphi_1$ oraz
$C_{k+1}=(C_k\land\varphi_{k+1})$ dla $1\le k<n$.
To jedynie konwencja zapisu, nie nowy spójnik.
\end{uwaga}

\section{Jak badać formuły? Wartościowania i tablice prawdy}
\label{sec:valuations}
\begin{definicja}[Wartościowanie]\label{def:valuation}
\emph{Wartościowaniem} dla rozważanej formuły nazywamy
przypisanie każdej występującej w niej literze zdaniowej
wartości $0$ lub $1$. Wartość całej formuły obliczamy od
najbardziej wewnętrznych jej części, korzystając z tablic
spójników logicznych. Ponieważ każde wystąpienie spójnika ma
argumenty będące krótszymi formułami, takie obliczanie kończy
się po skończonej liczbie kroków. Dwa wystąpienia tej samej litery
mają w jednym wartościowaniu tę samą wartość. Rozszerzoną
w ten sposób wartość formuły $\varphi$ oznaczamy $v(\varphi)$.
\end{definicja}

\begin{twierdzenie}[Jednoznaczność wartości formuły]
\label{thm:valuation-well-defined-xli}
Niech $\varphi$ będzie formułą zbudowaną zgodnie
z definicją~\ref{def:propositional-grammar-xli} i niech
$v$ przypisuje każdej literze występującej w $\varphi$
wartość $0$ albo $1$. Wówczas z tablic spójników można
wyznaczyć \emph{dokładnie jedną} wartość $v(\varphi)$.
\end{twierdzenie}
\begin{proof}
Rozpatrujemy kolejne etapy budowy formuły. W etapie
początkowym formuła jest pojedynczą literą; ma wartość
przyznaną jej przez $v$. Załóżmy, że wartości wszystkich
formuł powstałych w pierwszych $m$ etapach są jednoznaczne.
Jeśli w kolejnym etapie utworzono negację $\neg\alpha$,
to wartość $\alpha$ jest już znana, a tabela negacji
wskazuje dokładnie jedną wartość nowej formuły.
Jeżeli utworzono $(\alpha\circ\beta)$, gdzie $\circ$
jest jednym z czterech spójników dwuargumentowych,
to obie wartości $v(\alpha),v(\beta)$ są znane;
odpowiednia tabela wskazuje dokładnie jedną wartość wyniku.
Budowa $\varphi$ jest skończona, więc po skończonej
liczbie takich kroków otrzymamy jej wartość.
W pełnym zapisie pojedyncza litera jest rozpoznawana
jednoznacznie. Każda inna formuła zaczyna się od $\neg$
albo od nawiasu otwierającego. W pierwszym przypadku reszta
napisu stanowi argument negacji. W drugim spójnik dwuargumentowy
znajduje się wewnątrz zewnętrznej pary nawiasów, poza wszystkimi
wewnętrznymi parami; położenie tego spójnika wyznacza obie
podformuły. Stąd rozbiór formuły jest jednoznaczny.
Powtarzając to rozumowanie dla podformuł, widzimy, że żaden
krok obliczania wartości nie pozostawia dowolności.
\end{proof}

\begin{przyklad}[Obliczenie od wewnątrz]
Rozważmy $\varphi=(p\Rightarrow q)\land\neg q$.
Gdy $p=0$ i $q=0$, otrzymujemy kolejno:
$p\Rightarrow q=1$, $\neg q=1$, a zatem $\varphi=1$.
Gdy $p=1$ i $q=0$, mamy $p\Rightarrow q=0$ i
$\neg q=1$, więc $\varphi=0$. Zmiana jednego zdania
prostego może zmienić wartość całej formuły.
\end{przyklad}

\begin{definicja}[Tautologia, sprzeczność, spełnialność]\label{def:tautology}
Formułę nazywamy \emph{tautologią}, gdy przy \emph{każdym}
wartościowaniu jest prawdziwa. Jest \emph{sprzeczna}, gdy
przy każdym wartościowaniu jest fałszywa.
Jest \emph{spełnialna}, gdy istnieje \emph{przynajmniej jedno}
wartościowanie, w którym jest prawdziwa. Formuły
$\varphi$ i $\psi$ są \emph{logicznie równoważne},
gdy przy każdym wspólnym wartościowaniu mają taką samą
wartość logiczną; zapisujemy $\varphi\equiv\psi$.
Symbol $\equiv$ opisuje relację między formułami;
nie jest kolejnym spójnikiem tworzącym zdanie złożone.
Przy porównywaniu dwóch formuł wartościowanie obejmuje
\emph{wszystkie} litery występujące w którejkolwiek z nich,
także jeśli część liter występuje tylko w jednej formule.
\end{definicja}

\begin{przyklad}[Trzy różne sytuacje]
$p\lor\neg p$ jest tautologią, gdyż niezależnie od wartości
$p$ co najmniej jedno ze zdań $p,\neg p$ jest prawdziwe.
$p\land\neg p$ jest sprzecznością, gdyż nigdy nie mogą
być prawdziwe równocześnie. Formuła $p\land q$ jest
spełnialna (weźmy $p=q=1$), ale nie jest tautologią
(weźmy $p=1,q=0$).
\end{przyklad}

\begin{twierdzenie}[Test równoważności]\label{thm:equiv-test}
Dla dowolnych formuł zdaniowych $\varphi,\psi$ zachodzi:
$\varphi\equiv\psi$ wtedy i tylko wtedy, gdy
$\varphi\Leftrightarrow\psi$ jest tautologią.
\end{twierdzenie}
\begin{proof}
Ustalmy dowolne wartościowanie występujących w formułach
liter zdaniowych. Z definicji równoważności spójnikowej
$\varphi\Leftrightarrow\psi$ jest w tym wartościowaniu
prawdziwe dokładnie wtedy, gdy obie formuły mają tę samą
wartość. Zatem żądanie, aby działo się tak dla wszystkich
wartościowań, jest jednocześnie definicją tautologii
$\varphi\Leftrightarrow\psi$ i definicją relacji
$\varphi\equiv\psi$. Dowiedliśmy obu kierunków.
\end{proof}

\begin{twierdzenie}[Skończony test tablicowy]\label{thm:truth-table-decision-xli}
Jeżeli formuła $\varphi$ zawiera dokładnie $m$ różnych
liter zdaniowych, gdzie $m\ge1$, to ma dokładnie $2^m$
wartościowań. Zbadanie wszystkich tych wartościowań
pozwala rozstrzygnąć, czy $\varphi$ jest tautologią,
sprzecznością albo formułą spełnialną.
\end{twierdzenie}
\begin{proof}
Pierwszej literze możemy przypisać dwie wartości.
Po ustaleniu wartości pierwszych $k$ liter dla
następnej są znów dwa niezależne wybory. Zatem
liczba pełnych przypisań wynosi $2\cdot\ldots\cdot2=2^m$.
Każde przypisanie określa wartość $\varphi$ na mocy
twierdzenia~\ref{thm:valuation-well-defined-xli}.
Jeżeli każda z $2^m$ wartości jest równa $1$, definicja
daje tautologię. Jeśli każda jest równa $0$, otrzymujemy
sprzeczność. W pozostałym przypadku pojawia się co
najmniej jedna wartość $1$ i co najmniej jedna wartość $0$:
formuła jest spełnialna, ale nie jest tautologią.
\end{proof}

\begin{przyklad}[Trzy litery, lecz tylko osiem przypadków]
Rozważmy $(p\land q)\Rightarrow r$. Litery $p,q,r$
są trzy, więc wartościowań jest $2^3=8$. Implikacja
ma wartość $0$ tylko przy $p=q=1$, $r=0$;
w każdym z siedmiu pozostałych wierszy ma wartość $1$.
Jest zatem spełnialna, ale nie jest tautologią.
Powtórzone wystąpienia tej samej litery nie są
nowymi niezależnymi wyborami, np. $p\land p$
ma tylko dwa wartościowania, a nie cztery.
\end{przyklad}

\section{Podstawowe prawa rachunku zdań z dowodami}\label{sec:propositional-laws}
\begin{twierdzenie}[Prawa de Morgana]\label{thm:de-morgan-logic}
Dla dowolnych formuł $p,q$ zachodzą równoważności logiczne
\[
\neg(p\land q)\equiv(\neg p\lor\neg q),\qquad
\neg(p\lor q)\equiv(\neg p\land\neg q).
\]
\end{twierdzenie}
\begin{proof}
Pokażemy wszystkie cztery możliwe pary wartości formuł
$p$ i $q$. Jeżeli formuły mają wspólne litery, nie każda
para musi wystąpić przy pewnym wartościowaniu. Wystarcza jednak
sprawdzenie wszystkich czterech możliwości, bo każde
wartościowanie daje jedną z nich.
\[
\begin{array}{cc|cc|cc}
p&q&\neg(p\land q)&\neg p\lor\neg q&
\neg(p\lor q)&\neg p\land\neg q\\\hline
1&1&0&0&0&0\\
1&0&1&1&0&0\\
0&1&1&1&0&0\\
0&0&1&1&1&1
\end{array}
\]
Dwie kolumny opisujące pierwszą parę formuł są identyczne
w każdym wierszu; to samo dotyczy drugiej pary.
Z definicji równoważności logicznej otrzymujemy obie tezy.
\end{proof}

\begin{twierdzenie}[Rozwinięcie i przeciwstawienie implikacji]
\label{thm:implication-laws}
Dla dowolnych formuł $p,q$ zachodzą
\[
(p\Rightarrow q)\equiv(\neg p\lor q),\qquad
(p\Rightarrow q)\equiv(\neg q\Rightarrow\neg p).
\]
W szczególności $(p\Rightarrow q)\land p\Rightarrow q$
jest tautologią (\emph{modus ponens}).
\end{twierdzenie}
\begin{proof}
Rozpatrzmy cztery możliwe pary wartości formuł $p,q$;
nie wymagamy, aby każda z nich była realizowana przy pewnym
wspólnym wartościowaniu:
\[
\begin{array}{cc|ccc}
p&q&p\Rightarrow q&\neg p\lor q&\neg q\Rightarrow\neg p\\\hline
1&1&1&1&1\\
1&0&0&0&0\\
0&1&1&1&1\\
0&0&1&1&1
\end{array}
\]
W każdym wierszu trzy ostatnie wartości są równe, co
dowodzi obu równoważności. Aby sprawdzić ostatnią tezę,
rozpatrzmy dowolne wartościowanie. Jeśli koniunkcja
$(p\Rightarrow q)\land p$ jest fałszywa, całe zdanie
warunkowe jest prawdziwe z definicji implikacji.
Jeśli koniunkcja jest prawdziwa, to $p$ jest prawdziwe,
a ponieważ $p\Rightarrow q$ też jest prawdziwe, $q$
nie może być fałszywe: w przeciwnym razie implikacja
byłaby fałszywa. Zatem $q$ jest prawdziwe, więc
cała rozpatrywana implikacja jest prawdziwa również
i w tym przypadku.
\end{proof}

\begin{twierdzenie}[Negacja implikacji]\label{thm:negated-implication-xli}
Dla dowolnych formuł $p,q$ zachodzi
$\neg(p\Rightarrow q)\equiv p\land\neg q$.
\end{twierdzenie}
\begin{proof}
Ustalmy wartościowanie. Na mocy definicji implikacji
$p\Rightarrow q$ jest fałszywe wtedy i tylko wtedy,
gdy jednocześnie $p=1$ i $q=0$. Jego negacja jest
więc prawdziwa dokładnie przy $p=1$ i $q=0$.
Z drugiej strony $p\land\neg q$ jest prawdziwe
dokładnie wtedy, gdy $p=1$ oraz $\neg q=1$,
czyli $p=1$ i $q=0$. W obu przypadkach zbiór
wartościowań dających wartość $1$ jest taki sam,
a pozostałe wartościowania dają $0$. Otrzymaliśmy zatem
żądaną równoważność logiczną.
\end{proof}

\begin{przyklad}[Obliczenie kontrprzykładu do implikacji]
Niech $p$ oznacza „$12$ jest podzielne przez $4$”,
a $q$ --- „$12$ jest liczbą nieparzystą”. Wtedy
$p=1$, $q=0$, więc $p\Rightarrow q=0$ oraz
$p\land\neg q=1$. Zmienienie poprzednika na
zdanie fałszywe usunęłoby kontrprzykład do implikacji.
\end{przyklad}

\begin{twierdzenie}[Podwójna negacja i rozdzielność]
\label{thm:logic-distributivity}
Dla dowolnych formuł $p,q,r$ zachodzą
\begin{align*}
\neg\neg p&\equiv p,\\
p\land(q\lor r)&\equiv(p\land q)\lor(p\land r),\\
p\lor(q\land r)&\equiv(p\lor q)\land(p\lor r).
\end{align*}
\end{twierdzenie}
\begin{proof}
Dla podwójnej negacji wystarczy zauważyć, że przy
$p=1$ mamy $\neg p=0$ i $\neg\neg p=1$, zaś przy
$p=0$ mamy $\neg p=1$ i $\neg\neg p=0$.

Wykażmy pierwszą rozdzielność, analizując dwa przypadki.
Jeśli $p=0$, to lewa strona jest fałszywa, a prawa
również: oba składniki jej alternatywy zawierają
fałszywy czynnik $p$. Jeśli $p=1$, to lewa strona
przyjmuje dokładnie wartość $q\lor r$, a prawa strona
jest $(1\land q)\lor(1\land r)$, która także
przyjmuje wartość $q\lor r$. Oba przypadki wyczerpują
wszystkie wartościowania, ponieważ $p$ musi mieć wartość
$0$ lub $1$, niezależnie od wartości $q,r$.

Podobnie przy $p=1$ obie strony drugiej rozdzielności
są prawdziwe, a przy $p=0$ obie sprowadzają się do
$q\land r$. Stąd także druga równoważność.
\end{proof}

\begin{uwaga}[Dlaczego przykład liczbowy nie zastępuje dowodu?]
Sprawdzenie prawa logicznego na jednym przykładzie potwierdza
jedynie jedno wartościowanie. Aby wykazać tautologię,
trzeba zbadać \emph{wszystkie} wartościowania albo podać
ogólny argument obejmujący każdą możliwość. Z kolei do
\emph{obalenia} zdania głoszącego równoważność wystarczy
\emph{jedno} wartościowanie, przy którym formuły są różne.
\end{uwaga}

\section{Postacie normalne i zupełność funkcyjna}
\label{sec:lst-propositional-normal-forms}
Tablica prawdy pozwala odczytać wartości danej formuły.
Można również rozwiązać zadanie odwrotne: mając ustaloną
tablicę wartości, zbudować formułę, która ją realizuje.
W ten sposób otrzymamy systematyczną metodę zapisywania
i upraszczania formuł.

\begin{definicja}[Literały i postacie normalne]
\label{def:lst-dnf-cnf}
Literałem nazywamy literę zdaniową $p$ albo jej negację $\neg p$.
Formuła jest w alternatywnej postaci normalnej
(ang. \emph{disjunctive normal form}, DNF), gdy jest
alternatywą koniunkcji literałów. Jest w koniunkcyjnej
postaci normalnej (ang. \emph{conjunctive normal form}, CNF),
gdy jest koniunkcją alternatyw literałów.
Alternatywę literałów w CNF nazywamy klauzulą.
\end{definicja}
Pojedynczy literał jest zarówno koniunkcją, jak i alternatywą
jednoelementową. Dla $n\geq1$ puste alternatywy i koniunkcje
możemy zastępować odpowiednio przez $p_1\land\neg p_1$
oraz $p_1\lor\neg p_1$. Symbole $\bot,\top$ traktujemy
w tej części jako skróty tych formuł, a nie nowe symbole
w definicji składni z
Definicji~\ref{def:propositional-grammar-xli}.

\begin{twierdzenie}[Konstrukcja postaci normalnych z tablicy prawdy]
\label{thm:lst-canonical-normal-forms}
Niech $n\geq1$ i $f:\{0,1\}^n\to\{0,1\}$.
Dla $a=(a_1,\ldots,a_n)\in\{0,1\}^n$ określamy
\[
\begin{split}
 \ell_i^a&=\begin{cases}p_i,&a_i=1,\\\neg p_i,&a_i=0,\end{cases}\\
 m_a&=\bigwedge_{i=1}^n\ell_i^a,\qquad
 c_a=\bigvee_{i=1}^n\neg\ell_i^a.
\end{split}
\]
W wyrażeniu $c_a$ podwójną negację zastępujemy
odpowiednim literałem dodatnim.
Wówczas formuły
\[
 D_f=\bigvee_{a:\,f(a)=1}m_a,\qquad
 C_f=\bigwedge_{a:\,f(a)=0}c_a
\]
mają tablicę prawdy określoną przez $f$. W szczególności
każda formuła jest równoważna pewnej formule w DNF
oraz pewnej formule w CNF.
\end{twierdzenie}
\begin{proof}
Ustalmy wartościowanie $p_i=b_i$ dla
$b=(b_1,\ldots,b_n)\in\{0,1\}^n$.
Literał $\ell_i^a$ ma wartość $1$ dokładnie wtedy,
gdy $b_i=a_i$. Koniunkcja $m_a$ ma zatem wartość $1$
dokładnie dla $b=a$. Alternatywa wszystkich $m_a$,
dla których $f(a)=1$, ma wartość $1$ wtedy i tylko wtedy,
gdy $f(b)=1$. Otrzymujemy tablicę $D_f$.
Z kolei $c_a$ ma wartość $0$ dokładnie wtedy,
gdy wszystkie $\neg\ell_i^a$ mają wartość $0$,
czyli gdy $b=a$. Koniunkcja klauzul odpowiadających
wierszom, w których $f$ przyjmuje $0$, ma wartość $0$
dokładnie w tych wierszach. Otrzymujemy tablicę $C_f$.
W przypadku funkcji stałej obowiązuje podana wyżej
konwencja dotycząca pustych działań. Dla tablicy
dowolnej formuły podstawiamy jej funkcję wartości
za $f$ i stosujemy definicję równoważności logicznej.
\end{proof}
Wszystkie litery $p_1,\ldots,p_n$ występujące
w poszczególnych członach dają postacie pełne (kanoniczne).
Dla ustalonej kolejności liter i wierszy wybór członów
jest jednoznaczny, ponieważ każdy odpowiada dokładnie
jednemu wierszowi tablicy prawdy. Nie oznacza to
jednoznaczności dowolnej DNF lub CNF: równoważna
formuła może zawierać mniej literałów.

\begin{przyklad}[Dokładnie jedna z dwóch przesłanek]
\label{ex:lst-xor-normal-forms}
Niech $f(p,q)$ ma wartość $1$, gdy dokładnie jedna
z liter $p,q$ ma wartość $1$.
\[
\begin{array}{c|c|c|l|l}
 p&q&f&\text{człon DNF}&\text{klauzula CNF}\\ \hline
 0&0&0&-&p\lor q\\
 0&1&1&\neg p\land q&-\\
 1&0&1&p\land\neg q&-\\
 1&1&0&-&\neg p\lor\neg q
\end{array}
\]
Z wierszy prawdziwych odczytujemy
$(\neg p\land q)\lor(p\land\neg q)$,
a z fałszywych $(p\lor q)\land(\neg p\lor\neg q)$.
Obie formuły wyrażają tę samą funkcję alternatywy
rozłącznej. W drugim zapisie pierwszy nawias wymaga
co najmniej jednej prawdy, a drugi wyklucza dwie prawdy.
\end{przyklad}

\begin{przyklad}[Upraszczanie postaci pełnej]
\label{ex:lst-dnf-simplification}
Z dwóch sąsiednich wierszy tablicy, różniących się
tylko wartością $q$, możemy otrzymać krótszy człon:
\[
 (p\land q)\lor(p\land\neg q)
 \equiv p\land(q\lor\neg q)\equiv p.
\]
Pierwsza równoważność wynika z rozdzielności
(Twierdzenie~\ref{thm:logic-distributivity}),
a druga z faktu, że $q\lor\neg q$ przyjmuje zawsze
wartość $1$. Konstrukcja z tablicy daje więc
poprawną formułę, lecz nie gwarantuje najmniejszej
liczby spójników.
\end{przyklad}

\begin{definicja}[Zupełność funkcyjna spójników]
\label{def:lst-functional-completeness}
Rodzinę spójników nazywamy funkcyjnie zupełną, jeżeli
dla każdego $n\geq1$ i każdej funkcji
$f:\{0,1\}^n\to\{0,1\}$ istnieje formuła w literach
$p_1,\ldots,p_n$, używająca wyłącznie spójników
tej rodziny, której funkcją wartości jest $f$.
\end{definicja}

\begin{twierdzenie}\label{thm:lst-complete-connectives}
Funkcyjnie zupełne są rodziny
$\{\neg,\land,\lor\}$, $\{\neg,\land\}$,
$\{\neg,\lor\}$ i $\{\neg,\Rightarrow\}$.
\end{twierdzenie}
\begin{proof}
Pierwsza teza wynika z konstrukcji DNF
w Twierdzeniu~\ref{thm:lst-canonical-normal-forms}.
Alternatywę można zastąpić przez
$\neg(\neg p\land\neg q)$, a koniunkcję przez
$\neg(\neg p\lor\neg q)$, korzystając z praw
de Morgana (Twierdzenie~\ref{thm:de-morgan-logic}).
Dla rodziny $\{\neg,\Rightarrow\}$ mamy
$p\lor q\equiv(\neg p\Rightarrow q)$:
obie formuły są fałszywe dokładnie dla $p=q=0$.
Stąd możemy zbudować alternatywę, negację i,
przez prawo de Morgana, koniunkcję.
Sukcesywne zastępowanie spójników w dowolnej DNF
zachowuje wartość całej formuły, co wynika
z rekurencyjnej definicji wartościowania.
\end{proof}

\begin{twierdzenie}[Pojedynczy spójnik wystarcza]
\label{thm:lst-nand-nor}
Każda z jednoelementowych rodzin zawierających
spójnik NAND (kreska Sheffera) lub NOR (strzałka Peirce'a)
jest funkcyjnie zupełna.
\end{twierdzenie}
\begin{proof}
Piszemy $p\uparrow q=\neg(p\land q)$ oraz
$p\downarrow q=\neg(p\lor q)$. Dla NAND zachodzi
\[
 \neg p\equiv p\uparrow p,\qquad
 p\land q\equiv(p\uparrow q)\uparrow(p\uparrow q).
\]
Dla NOR zachodzi
\[
 \neg p\equiv p\downarrow p,\qquad
 p\lor q\equiv(p\downarrow q)\downarrow(p\downarrow q).
\]
Wzory wynikają bezpośrednio z idempotentności
koniunkcji i alternatywy oraz prawa podwójnej negacji;
można też sprawdzić cztery wiersze tablicy.
Twierdzenie~\ref{thm:lst-complete-connectives}
daje zupełność obu rodzin.
\end{proof}

\begin{stwierdzenie}[Rodziny niezupełne funkcyjnie]
\label{prop:lst-incomplete-connectives}
Rodzina $\{\land,\lor\}$ nie jest funkcyjnie zupełna,
nawet po dołączeniu obu stałych logicznych.
Sama implikacja bez stałej fałszu również nie jest
funkcyjnie zupełna.
\end{stwierdzenie}
\begin{proof}
Na $\{0,1\}^n$ przyjmujemy porządek współrzędny,
$0<1$. Rzuty i funkcje stałe są niemalejące,
a minimum oraz maksimum dwóch funkcji niemalejących
są niemalejące. Indukcja po budowie formuły
pokazuje, że każda formuła z koniunkcji,
alternatywy i stałych wyznacza funkcję niemalejącą.
Negacja nie jest niemalejąca, ponieważ jej wartości
w $0$ i $1$ to odpowiednio $1$ i $0$.
Dla samej implikacji indukcja pokazuje, że przy
wartościowaniu wszystkich liter przez $1$ każda
formuła ma wartość $1$: w kroku indukcyjnym
korzystamy z $1\Rightarrow1=1$. Nie można więc
zbudować negacji ani funkcji stale równej $0$.
\end{proof}

\begin{wniosek}[Ile istnieje funkcji prawdziwościowych?]
\label{cor:lst-boolean-function-count}
Dla ustalonej listy $n\geq1$ liter zdaniowych istnieje
dokładnie $2^{2^n}$ klas równoważności logicznej formuł
w tych literach.
\end{wniosek}
\begin{proof}
Tablica ma $2^n$ wierszy, w każdym można niezależnie
wybrać $0$ albo $1$. Jest zatem $2^{2^n}$ tablic.
Każdą realizuje formuła z
Twierdzenia~\ref{thm:lst-canonical-normal-forms}.
Dwie formuły są równoważne dokładnie wtedy, gdy
mają tę samą tablicę prawdy. Stąd podana liczba klas.
\end{proof}
Przykładowo dla dwóch liter istnieje $16$ funkcji
prawdziwościowych, chociaż samych formuł jest
nieskończenie wiele. Zupełność funkcyjna dotyczy
reprezentowania funkcji prawdziwościowych; nie jest
tym samym co pełność rachunku dowodów, która dotyczy
przejścia od konsekwencji semantycznej do dowodliwości.

\section{Przesłanki, wnioski i kontrprzykłady}\label{sec:logical-consequence}
\begin{definicja}[Poprawność wnioskowania zdaniowego]
\label{def:consequence}
Mówimy, że formuła $\psi$ \emph{wynika logicznie}
z przesłanek $\varphi_1,\ldots,\varphi_n$, gdy nie istnieje
\emph{wspólne} wartościowanie wszystkich liter występujących
w przesłankach lub we wniosku, przy którym wszystkie przesłanki
są prawdziwe, a $\psi$ jest fałszywa. Zapisujemy
$\varphi_1,\ldots,\varphi_n\models\psi$.
Przy braku przesłanek $\models\psi$ oznacza, że
$\psi$ jest tautologią.
\end{definicja}

\begin{twierdzenie}[Kryterium poprawności wnioskowania]
\label{thm:inference-test}
Niech $\varphi_1,\ldots,\varphi_n,\psi$ będą formułami
zdaniowymi, gdzie $n\ge1$. Następujące warunki są
równoważne:
\begin{enumerate}[label=(\roman*)]
\item $\varphi_1,\ldots,\varphi_n\models\psi$;
\item $((\varphi_1\land\cdots\land\varphi_n)\Rightarrow\psi)$
      jest tautologią.
\end{enumerate}
\end{twierdzenie}
\begin{proof}
Ustalmy dowolne wartościowanie. Implikacja z punktu
(ii) jest fałszywa dokładnie wtedy, gdy jej
poprzednik jest prawdziwy, a następnik fałszywy.
Poprzednik jest koniunkcją wszystkich przesłanek,
więc jest prawdziwy dokładnie wtedy, gdy wszystkie
przesłanki są prawdziwe. Zatem istnienie
wartościowania obalającego tautologiczność w (ii)
jest równoważne istnieniu wartościowania obalającego
poprawność w (i). Zaprzeczenie obu tych stwierdzeń
daje żądaną równoważność.
\end{proof}

\begin{twierdzenie}[Jednolite kryterium wynikania dla skończonej listy przesłanek]
\label{thm:empty-premises-consequence-lxiii}
Niech $n\in\{0,1,2,\ldots\}$, niech $\varphi_1,\ldots,\varphi_n$ będą formułami
rachunku zdań, a $\psi$ dowolną formułą. Dla $n=0$ przyjmijmy, że
koniunkcja pustej listy ma stałą wartość logiczną $1$ (oznaczamy ją
metajęzykowo przez $\top$). Wówczas
\[
\varphi_1,\ldots,\varphi_n\models\psi
\quad\Longleftrightarrow\quad
\models\bigl((\varphi_1\land\cdots\land\varphi_n)\Rightarrow\psi\bigr),
\]
gdzie przy $n=0$ prawą stronę rozumiemy jako $\models(\top\Rightarrow\psi)$.
Symbol $\top$ jest tutaj skrótem semantycznym, a nie dodatkową
literą w gramatyce z definicji~\ref{def:propositional-grammar-xli}.
\end{twierdzenie}
\begin{proof}
Dla $n\ge1$ jest to twierdzenie~\ref{thm:inference-test}.
Dla $n=0$ definicja~\ref{def:consequence} mówi, że
$\models\psi$ wtedy i tylko wtedy, gdy każde wartościowanie
nadaje $\psi$ wartość $1$. Ponieważ $\top$ ma zawsze wartość $1$,
implikacja $\top\Rightarrow\psi$ ma przy każdym wartościowaniu
dokładnie tę samą wartość co $\psi$. Stąd równoważność również
dla pustej listy. Nie trzeba wybierać żadnego wartościowania
spełniającego przesłanki: przesłanek po prostu nie ma.
\end{proof}
\begin{przyklad}[Pusta lista nie jest listą sprzeczną]
Dla $\psi=p\lor\neg p$ zarówno $\psi$, jak i
$\top\Rightarrow\psi$ mają wartość $1$ przy $p=0$
i przy $p=1$. Dlatego oba stwierdzenia w kryterium są
prawdziwe. Dla $\psi=p$ oba stwierdzenia o wszystkich
wartościowaniach są fałszywe: $p=0$ daje kontrprzykład.
Natomiast dla listy $(p,\neg p)$
wynikanie dowolnej formuły jest prawdziwe z innego powodu:
żadne wartościowanie nie spełnia wszystkich przesłanek.
\end{przyklad}

\begin{uwaga}[Przesłanki sprzeczne i brak przesłanek]
\label{rem:explosion-semantic-xli}
Nie należy mylić prawdziwości przesłanek
z poprawnością reguły wnioskowania. Dla dowolnych formuł
$p,q$ zachodzi $p,\neg p\models q$. Rzeczywiście,
żadne wartościowanie nie czyni obu przesłanek prawdziwymi:
jeżeli $p=1$, to $\neg p=0$, a jeżeli $p=0$, to pierwsza
przesłanka jest fałszywa. Nie istnieje więc wartościowanie
z prawdziwymi wszystkimi przesłankami i fałszywym wnioskiem.
Jest to semantyczna postać \emph{zasady eksplozji}; nie znaczy,
że dowolne zdanie $q$ jest prawdziwe. Podobnie
$\models p\lor\neg p$, lecz nie zachodzi
$\models p$ dla samej litery $p$, bo wartościowanie
$p=0$ jest kontrprzykładem.
\end{uwaga}

\begin{przyklad}[Niepoprawne wnioskowanie: potwierdzanie następnika]
Rozumowanie „$p\Rightarrow q$; $q$; zatem $p$” jest
niepoprawne. Wystarczy przyjąć $p=0$, $q=1$:
wtedy obie przesłanki są prawdziwe, ale wniosek
$p$ jest fałszywy. Konkretny przykład: „jeżeli
liczba jest podzielna przez $4$, to jest parzysta;
liczba $6$ jest parzysta; zatem jest podzielna przez $4$”.
\end{przyklad}

\begin{przyklad}[Niepoprawne wnioskowanie: zaprzeczanie poprzednikowi]
Z $p\Rightarrow q$ oraz $\neg p$ nie wynika $\neg q$.
Przyjmując $p=0$, $q=1$, otrzymujemy prawdziwe
przesłanki i fałszywy wniosek. Jest to ten sam
błąd co twierdzenie, że liczba niepodzielna przez $4$
nie może być parzysta.
\end{przyklad}

\begin{twierdzenie}[Trzy poprawne reguły wnioskowania]
\label{thm:valid-arguments}
Dla dowolnych formuł $p,q,r$ zachodzą
\begin{align*}
p\Rightarrow q,\quad\neg q&\models\neg p&&\text{(modus tollens)},\\
p\Rightarrow q,\quad q\Rightarrow r&\models p\Rightarrow r
   &&\text{(przechodniość implikacji)},\\
p\lor q,\quad\neg p&\models q&&\text{(sylogizm rozłączny)}.
\end{align*}
\end{twierdzenie}
\begin{proof}
W pierwszym przypadku przyjmijmy, że obie przesłanki
są prawdziwe. Ponieważ $\neg q$ jest prawdziwe,
$q$ jest fałszywe. Gdyby $p$ było prawdziwe,
$p\Rightarrow q$ byłoby fałszywe, sprzecznie z
założeniem. Zatem $p$ jest fałszywe, czyli
$\neg p$ jest prawdziwe.

W drugim przypadku załóżmy prawdziwość obu implikacji.
Chcemy uzasadnić $p\Rightarrow r$. Gdy $p$ jest
fałszywe, implikacja ta jest prawdziwa z definicji.
Gdy $p$ jest prawdziwe, pierwsza przesłanka wymusza
prawdziwość $q$, a druga prawdziwość $r$.
Wtedy również $p\Rightarrow r$ jest prawdziwa.

W trzecim przypadku $p\lor q$ jest prawdziwe, więc
co najmniej jedna z wartości $p,q$ wynosi $1$.
Jednocześnie $\neg p$ jest prawdziwe, a więc $p=0$.
Pozostaje $q=1$, co jest dokładnie tezą.
\end{proof}

\begin{przyklad}[Jak odczytywać wzory, aby nie mylić spójnika z wnioskowaniem]
Zapis $(p\Rightarrow q)$ jest \emph{formułą}: można obliczyć
jego wartość z tablicy prawdy. Natomiast zapis
$p\models q$ \emph{nie jest formułą rachunku zdań}, lecz
metamatematycznym stwierdzeniem, że nie istnieje wartościowanie
z $p=1$ i $q=0$. W przypadku jednej przesłanki zachodzi
$p\models q$ wtedy i tylko wtedy, gdy $p\Rightarrow q$ jest
tautologią. Podobnie $p\equiv q$ opisuje zgodność wartości
formuł przy wszystkich wartościowaniach, a
$p\Leftrightarrow q$ jest formułą, której wartość badamy.
\end{przyklad}

\section{Zbiory przesłanek i konsekwencja logiczna}
\label{sec:lst-propositional-consequence}
Dotychczas rozważaliśmy skończone listy przesłanek.
Możemy jednak ustalić także zbiór nieskończony, na przykład
nieskończoną listę warunków, które mają zachodzić jednocześnie.
Nadal każda pojedyncza formuła ma skończoną budowę.

\begin{definicja}[Spełnialność zbioru formuł i wynikanie]
\label{def:lst-propositional-theory}
Niech $\Gamma$ będzie dowolnym zbiorem formuł zdaniowych.
Wartościowanie $v$ spełnia $\Gamma$, gdy $v(\gamma)=1$
dla każdego $\gamma\in\Gamma$. Zbiór jest spełnialny,
gdy istnieje takie wartościowanie.
Piszemy $\Gamma\models\varphi$, gdy każde wartościowanie
spełniające $\Gamma$ nadaje formule $\varphi$ wartość $1$.
\end{definicja}
Definicja rozszerza Definicję~\ref{def:consequence}.
Zbiór pusty jest spełnialny: spełnia go każde wartościowanie.
Nie wystarcza natomiast, by każda formuła z $\Gamma$
była z osobna spełnialna. Na przykład $p$ oraz $\neg p$
są z osobna spełnialne, ale nie mają wspólnego wartościowania
spełniającego obie formuły.

\begin{twierdzenie}[Kryterium przez niespełnialność]
\label{thm:lst-propositional-refutation}
Dla dowolnego zbioru przesłanek $\Gamma$ i formuł
$\varphi,\psi$ zachodzą równoważności
\[
\begin{split}
 \Gamma\models\varphi
 &\iff \Gamma\cup\{\neg\varphi\}
          \text{ jest niespełnialny},\\
 \Gamma\cup\{\varphi\}\models\psi
 &\iff\Gamma\models(\varphi\Rightarrow\psi).
\end{split}
\]
\end{twierdzenie}
\begin{proof}
W pierwszym przypadku zaprzeczenie lewej strony oznacza
istnienie wartościowania spełniającego $\Gamma$ i nadającego
$\varphi$ wartość $0$, czyli nadającego $\neg\varphi$ wartość $1$.
Jest to dokładnie spełnialność zbioru po prawej stronie.
W drugim przypadku, przy wartościowaniu spełniającym $\Gamma$,
implikacja $\varphi\Rightarrow\psi$ jest fałszywa dokładnie
wtedy, gdy $\varphi$ jest prawdziwa, a $\psi$ fałszywa.
Takie wartościowanie obala obie strony równoważności.
\end{proof}
Drugą równoważność nazywamy semantyczną postacią twierdzenia
o dedukcji. Nie zakładamy tu żadnego formalnego rachunku
dowodów. Syntaktyczne twierdzenie o dedukcji dla dedukcji
naturalnej omawiamy w teorii modeli.

\begin{stwierdzenie}[Własności konsekwencji]
\label{prop:lst-propositional-consequence-properties}
Dla zbiorów formuł $\Gamma,\Delta$ i formuły $\varphi$
zachodzą następujące własności:
\begin{enumerate}
 \item jeśli $\varphi\in\Gamma$, to $\Gamma\models\varphi$;
 \item jeśli $\Gamma\subseteq\Delta$ i $\Gamma\models\varphi$,
       to $\Delta\models\varphi$;
 \item jeśli $\Gamma\models\delta$ dla każdego $\delta\in\Delta$
       i $\Delta\models\varphi$, to $\Gamma\models\varphi$.
\end{enumerate}
\end{stwierdzenie}
\begin{proof}
Każde wartościowanie spełniające $\Gamma$ spełnia w szczególności
każdą należącą do niego formułę, co daje pierwszy punkt.
Jeśli $\Gamma\subseteq\Delta$, to każde wartościowanie spełniające
$\Delta$ spełnia także $\Gamma$. Daje to drugi punkt.
W trzecim, dowolne wartościowanie spełniające $\Gamma$
spełnia wszystkie formuły z $\Delta$, a więc także $\varphi$.
Argumenty pozostają ważne dla zbiorów niespełnialnych:
wówczas nie istnieje wartościowanie obalające konsekwencję.
\end{proof}
Druga własność nazywa się monotonicznością: dodanie przesłanek
nie unieważnia żadnego dotychczasowego wniosku klasycznego.
Trzecia pozwala zastąpić cały łańcuch pośrednich wnioskowań
jednym wynikaniem od przesłanek początkowych do końcowego wniosku.

\begin{twierdzenie}[Jednoczesne podstawianie formuł]
\label{thm:lst-propositional-substitution}
Niech $\sigma$ przyporządkowuje każdej literze zdaniowej $p$
formułę $\sigma(p)$. Rozszerzamy $\sigma$ na wszystkie formuły,
zachowując ich spójniki i podstawiając jednocześnie
przypisane formuły w miejsca liter. Jeżeli
$\Gamma\models\varphi$, to $\sigma(\Gamma)\models\sigma(\varphi)$,
gdzie $\sigma(\Gamma)=\{\sigma(\gamma):\gamma\in\Gamma\}$.
\end{twierdzenie}
\begin{proof}
Ustalmy dowolne wartościowanie $v$ i określmy wartościowanie
liter $w(p)=v(\sigma(p))$. Indukcja po budowie formuły $\theta$
pokazuje, że $v(\sigma(\theta))=w(\theta)$.
Dla litery wynika to z definicji $w$, a dla spójnika
obie strony powstają przez zastosowanie tej samej funkcji
prawdziwościowej do równych wartości argumentów.
Jeżeli $v$ spełnia $\sigma(\Gamma)$, to $w$ spełnia $\Gamma$.
Z założenia $w(\varphi)=1$, zatem $v(\sigma(\varphi))=1$.
\end{proof}
W szczególności podstawienie dowolnych formuł w miejsce liter
tautologii daje tautologię. Na przykład do tautologii
$p\lor\neg p$ możemy podstawić formułę $q\land r$ za $p$
i otrzymać $(q\land r)\lor\neg(q\land r)$.
Nie podstawiamy następnie rekurencyjnie za litery pojawiające się
we wstawionej formule: podstawienie jest jednoczesne.

\section{Skończoność przesłanek}
\label{sec:lst-propositional-compactness}
Wynikanie z nieskończonej liczby warunków może wyglądać na
istotnie silniejsze od wynikania ze skończonej listy.
W klasycznym rachunku zdań zachodzi jednak następujące
twierdzenie. Jego dowód podamy dla ustalonego wcześniej
przeliczalnego języka liter zdaniowych.

\begin{twierdzenie}[Zwartość rachunku zdań]
\label{thm:lst-propositional-compactness}
Zbiór formuł $\Gamma$ w języku o przeliczalnym zbiorze liter
zdaniowych jest spełnialny wtedy i tylko wtedy,
gdy każdy jego skończony podzbiór jest spełnialny.
\end{twierdzenie}
\begin{proof}
Jeżeli $v$ spełnia $\Gamma$, to spełnia każdy jego podzbiór.
Dowodzimy implikacji odwrotnej.
W przypadku skończonego zbioru liter dołączamy w razie potrzeby
przeliczalnie wiele nowych liter niewystępujących w $\Gamma$,
co nie zmienia spełnialności.
Oznaczmy litery przez $p_1,p_2,\ldots$.
Dla skończonego przypisania $a_1,\ldots,a_n\in\{0,1\}$
rozpatrujemy własność: każdy skończony podzbiór $\Gamma$
ma wartościowanie spełniające go i przyjmujące
wartości $a_i$ na literach $p_i$ dla $i\leq n$.

Przypisanie puste ma tę własność z założenia.
Jeśli ma ją przypisanie długości $n$, to co najmniej jedno
z jego dwóch rozszerzeń przez $a_{n+1}=0$ albo $a_{n+1}=1$
także ma tę własność.
Gdyby nie miało jej żadne, istniałyby dwa skończone podzbiory
$\Gamma_0,\Gamma_1\subseteq\Gamma$: pierwszy nie miałby
wartościowania zgodnego z rozszerzeniem przez $0$,
a drugi zgodnego z rozszerzeniem przez $1$.
Skończony zbiór $\Gamma_0\cup\Gamma_1$ nie miałby wtedy
wartościowania zgodnego z wyjściowym przypisaniem,
ponieważ dowolne wartościowanie musi nadać $p_{n+1}$
jedną z tych dwóch wartości. Otrzymujemy sprzeczność.

Budujemy więc ciąg $a_1,a_2,\ldots$, na każdym kroku
wybierając rozszerzenie przez $0$, jeśli zachowuje własność,
a w przeciwnym razie przez $1$. Definiujemy $v(p_i)=a_i$.
Dla dowolnej $\gamma\in\Gamma$ istnieje $n$ takie, że
wszystkie jej litery należą do $p_1,\ldots,p_n$,
gdyż formuła jest skończona.
Własność z etapu $n$ zastosowana do zbioru $\{\gamma\}$
daje wartościowanie $w$ spełniające $\gamma$ i zgodne
z $v$ na tych literach. Z rekurencyjnej definicji
wartościowania wynika $v(\gamma)=w(\gamma)=1$.
Tak jest dla każdej $\gamma$, zatem $v$ spełnia $\Gamma$.
\end{proof}

\begin{wniosek}[Skończony zbiór istotnych przesłanek]
\label{cor:lst-propositional-finite-support}
W przeliczalnym języku zdaniowym $\Gamma\models\varphi$
wtedy i tylko wtedy, gdy istnieje skończony podzbiór
$\Gamma_0\subseteq\Gamma$ taki, że $\Gamma_0\models\varphi$.
\end{wniosek}
\begin{proof}
Jeżeli $\Gamma\models\varphi$, to
$\Gamma\cup\{\neg\varphi\}$ jest niespełnialny na mocy
Twierdzenia~\ref{thm:lst-propositional-refutation}.
Ze zwartości wynika, że pewien jego skończony podzbiór
$\Sigma$ jest niespełnialny. Niech $\Gamma_0=\Sigma\cap\Gamma$.
Mamy $\Sigma\subseteq\Gamma_0\cup\{\neg\varphi\}$,
więc ten ostatni zbiór jest również niespełnialny.
Ponownie stosujemy kryterium przez niespełnialność.
Kierunek odwrotny wynika z monotoniczności
(Stwierdzenie~\ref{prop:lst-propositional-consequence-properties}).
\end{proof}

\begin{przyklad}[Nieskończony łańcuch implikacji]
\label{ex:lst-propositional-chain}
Rozważmy $\Gamma=\{p_1\}\cup\{p_i\Rightarrow p_{i+1}:i\geq1\}$.
Dla każdego ustalonego $n$ mamy $\Gamma\models p_n$:
od $p_1$ przechodzimy kolejno do $p_2,\ldots,p_n$.
Wystarcza skończony podzbiór złożony z $p_1$ i pierwszych
$n-1$ implikacji. Nie istnieje jednak jeden skończony
podzbiór $\Gamma_0$ wystarczający do wyprowadzenia wszystkich $p_n$.
Wybierzmy $N\geq1$ tak, by wszystkie litery występujące
w $\Gamma_0$ miały indeks co najwyżej $N$.
Wartościowanie nadające $p_1,\ldots,p_N$ wartość $1$
i pozostałym literom wartość $0$ spełnia $\Gamma_0$,
lecz nie spełnia $p_{N+1}$.
\end{przyklad}
Zwartość nie daje algorytmu rozstrzygającego spełnialność
dowolnie zadanego nieskończonego zbioru formuł.
Mówi o istnieniu skończonej przeszkody w razie niespełnialności,
nie o sposobie jej znalezienia przy dowolnej reprezentacji danych.
Zwartość logiki pierwszego rzędu ma osobny dowód
w teorii modeli.

\section{Czego jeszcze nie potrafi rachunek zdań?}\label{sec:why-quantifiers}
Niech $P(x)$ oznacza wypowiedź „liczba $x$ jest parzysta”.
Możemy sprawdzić $P(2)$ albo $P(3)$, ale sama logika zdań
nie ma jeszcze symbolu pozwalającego zbudować wypowiedź
„każda liczba ma pewną własność” lub „istnieje liczba
o tej własności”. Nie odróżnia też jednoznacznie wypowiedzi
„każdy ma swojego przyjaciela” od „istnieje ktoś,
kto jest przyjacielem wszystkich”. Do tego potrzebujemy
\emph{kwantyfikatorów}; wprowadzimy je w następnym rozdziale.

\begin{uwaga}[Pierwsze trzy rozróżnienia, które należy zapamiętać]
\emph{Zdanie prawdziwe} nie musi być \emph{tautologią}.
\emph{Implikacja} na ogół nie jest logicznie równoważna
swojej implikacji odwrotnej.
\emph{Poprawne wnioskowanie} nie zakłada samo przez się,
że jego przesłanki są faktycznie prawdziwe.
\end{uwaga}
